\documentclass[12pt]{article}
%^ Start of document
\usepackage{times}  %Makes text TimesNewRoman
\usepackage{setspace} %Allows you to set single or double space 
\usepackage{biblatex} %Imports biblatex package
\usepackage{graphicx} % Required for inserting images
\usepackage{authblk}
\usepackage[colorlinks=true, urlcolor=blue, linkcolor=blue]{hyperref}
\usepackage{subfig}
\usepackage{float}
\usepackage{tabularx}
\usepackage{multirow}
\usepackage[paperheight=11in,
   paperwidth=8.5in,
   top=20mm,
   bottom=20mm,
   left=40mm,
   right=40mm]{geometry}
\usepackage{moresize}
\usepackage{tikz}
\usepackage{xcolor}
\usepackage{physics}
\usepackage{amssymb}
\usepackage{mathrsfs}
\usepackage{amsmath}
\usepackage[document]{ragged2e}
\usepackage{comment}
\usepackage{xcolor}
\addbibresource{Qhw.bib}
\begin{document}
\singlespacing  %Sets single spacing for whole document
\title{Interdisciplinary Evolution and Research in Modern Physics and Biosciences}

\author[1]{Zachary Tullier}
\affil[1]{Student\\Physics Department\\ University of Houston - Clear Lake \\Houston, Texas\\U.S}
\date{}
\maketitle

\section*{Abstract}
This paper summarizes a wide-ranging presentation that discusses the evolution of physics and science education in response to interdisciplinary demands, especially in the post-pandemic world. The speaker outlines the shift in academic focus from traditional, discipline-specific research to problem-oriented and interdisciplinary approaches---particularly involving physics, biology, computational modeling, and medical science. Personal anecdotes, pedagogical philosophy, and ongoing student-driven projects (such as bacterial response to sound and fields) illustrate the deeply integrative nature of modern science.

\section*{Introduction}
The presentation opens by highlighting a transformation in scientific inquiry and education, where traditional boundaries have dissolved. Once dominated by siloed disciplines such as astrophysics or particle physics, today's scientific environment demands cross-disciplinary fluency---particularly as societal needs and job markets shift. The speaker underscores how events like the COVID-19 pandemic highlighted the essential nature of food and healthcare, driving academic and funding focus toward bioscience, medical applications, and sustainable energy.

\section*{Interdisciplinary Science and Career Implications}
Despite a personal passion for cosmology and particle physics, the speaker emphasizes the professional necessity for students to remain agile. Careers in theoretical fields like string theory may require students to pivot into applied domains such as financial modeling, biophysics, or systems biology. Computational skills, rather than field-specific content knowledge alone, often become the transferable asset in such transitions.

The importance of multi-particle systems, as opposed to single-particle basic science, is also stressed. The former finds relevance in materials science, biological systems, and astronomical structures, reinforcing the need for physicists to adapt their knowledge into broader contexts.

\section*{Physics as a Central Discipline}
Physics is championed as the core science that connects and enables multiple domains. Concepts from abstract mathematics, like differential geometry, find new utility in modeling molecular structures in biology. According to the speaker, a physicist trained in many-body physics can effectively engage with diverse systems---from stellar interiors to the human microbiome---due to shared structural and dynamical principles.

The talk includes a philosophical remark: physical systems behave independently of who observes them. Thus, the discipline used to interrogate these systems dictates the nature and depth of insight, not the system itself.

\section*{Departmental Updates and Student Guidance}
The speaker provides an honest evaluation of current faculty, highlighting individuals with strong research productivity---especially those leveraging computational tools. Experimental work, while more constrained due to equipment and time, is still deeply valued. Students are encouraged to consider both the nature of their interests and the practicalities of mentorship, research style, and communication fit when choosing advisors or research topics.

A notable point is that formalities (like having a committee chair from physics) do not prevent interdisciplinary collaboration, especially across departments like mechanical engineering or computer science.

\section*{Case Study: Bacterial Response to Environmental Perturbations}
A major portion of the latter half of the talk discusses an ongoing undergraduate-led research project investigating bacterial behavior (specifically \textit{E. coli} and \textit{Staphylococcus aureus}) under external influences such as sound waves and magnetic fields.

\begin{itemize}
    \item Differences between gram-positive and gram-negative bacteria and their membrane structures.
    \item The motility of bacteria in ionic fluid media and how external fields may impact bacterial dynamics more indirectly via the surrounding ions.
    \item Analogies drawn between nuclear decay equations and bacterial population growth models.
    \item A speculative but promising line of inquiry into how bacteria ``respond'' to vibrational stimuli (e.g., music), suggesting that mechanical waves might influence morphology or replication rates.
\end{itemize}

The project is presented as an example of how a physics-driven, model-building approach can yield insights in a completely foreign domain like microbiology.

\section*{Closing Thoughts}
The speaker encourages students to follow their scientific curiosity but also to remain grounded in pragmatic considerations---interests, mentors, job prospects, and communication styles. Passion is deemed a critical driver that allows students to overcome bureaucratic and academic hurdles.

\section*{Conclusion}
This presentation paints a compelling picture of modern scientific training as necessarily interdisciplinary, adaptive, and driven by both global needs and personal inquiry. Physics serves not only as a foundational discipline but also as a springboard into domains like medical physics, computational biology, and energy systems. Through storytelling and case studies, the speaker conveys an inclusive and intellectually rigorous vision for science education that transcends traditional silos and embraces complexity.

\end{document}
